% ---------------------------------------------------------------------------
% Standalone construct section, meant to be % ---------------------------------------------------------------------------
% Standalone construct section, meant to be % ---------------------------------------------------------------------------
% Standalone construct section, meant to be % ---------------------------------------------------------------------------
% Standalone construct section, meant to be \input{section_construct} from
% ssrn_diffusion_paper.tex so it can be revised independently of the rest.
%
% Depends on the host preamble for: \dkl{} (defined as $\Delta$KL), natbib,
% booktabs, amsmath. Citation keys used here all exist in the host
% bibliography: Murdock2017, Barron2018, vonGraevenitz2022.
%
% Naming decided 2026-08-06 after consultation. The signed axis is LEAD, and
% filings are LEADING or LAGGING; "currency" was rejected because readers
% import monetary logic into what is a claim about lexical surprise, and
% "resonance" because it names a mechanism this design cannot observe.
% ---------------------------------------------------------------------------

\section{The construct: how surprising a description is, and when}\label{sec:construct}

\subsection{Two questions asked of every filing}

The measure asks one thing twice. Take a filing's goods/services description,
and group every filing by its Nice class --- one of the 45 international goods
and services categories --- and by its year. Then ask:

\begin{enumerate}[noitemsep]
  \item How surprising would this wording have been to a reader who knew only
        what that industry had filed over the previous five years?
  \item How surprising does the same wording look to a reader who also knows
        what the industry filed over the following five years?
\end{enumerate}

\noindent The two answers are rarely the same, and the gap between them is
informative. Wording that was surprising when filed and unremarkable five
years later was surprising because it came first: the industry moved toward
it. Wording that read as ordinary when filed and unusual afterwards was behind
its industry rather than ahead of it, and was left behind as the field moved
on. Filings of the first kind are \emph{leading}; of the second,
\emph{lagging}.

Both quantities are computable for every filing in the corpus, including the
millions from firms that never patent, and neither requires knowing anything
about the firm.

\subsection{Antecedents, and what does not carry over}

The apparatus is that of \citet{Murdock2017}, developed on Darwin's reading
notebooks, and \citet{Barron2018}, who extended it to French Revolution
parliamentary debates. They score each document against the distribution of
documents before it and again against the distribution after it, and call the
two quantities \emph{novelty} and \emph{transience} and their signed difference
\emph{resonance}.

The arithmetic transfers exactly. The interpretation does not, in two ways that
shape everything downstream.

First, there is no single authoring agent. In a reading notebook the author of
the document is the person whose novelty is being measured, and the reference
stream is that author's own intellectual environment; a high value therefore
says something about a mind. Here each filing is scored against the stream of
\emph{other} firms' filings, and the author is typically counsel drafting for
scope of protection. The same number can only say that a description sits far
from what its industry was writing.

Second, \emph{resonance} names a mechanism this design cannot observe. A filing
can sit where its class was heading because it was imitated, because it and the
class responded to the same external development, or by coincidence, and
nothing here separates those. The vocabulary used below is therefore positional
rather than causal: a leading filing was ahead of its industry's language,
which implies the existence of a trajectory without asserting that the filing
caused it.

The nearest measurement precedent on this corpus is \citet{vonGraevenitz2022},
who flag goods/services tokens that are both new and fast-spreading and use
their geography to study diffusion. Their unit is the token and their novelty
flag is binary, assigned once a token's path is known; the measure here is
continuous and attaches to the filing, so every application carries a position
and that position can be crossed with the mark's own prosecution outcomes.

\subsection{Definitions}

Let $P_t$ be the topic distribution of a filing made in year $t$, and let
$Q_{\mathrm{past}}$ and $Q_{\mathrm{future}}$ aggregate same-class filings over
the $W = 5$ years before and after it. Write
\[
  K^{-} \;=\; \mathrm{KL}\!\left(P_t \,\|\, Q_{\mathrm{past}}\right),
  \qquad
  K^{+} \;=\; \mathrm{KL}\!\left(P_t \,\|\, Q_{\mathrm{future}}\right).
\]
Kullback--Leibler divergence is large when the filing puts weight where its
industry's filings put little, so $K^{-}$ is surprise measured against the
past and $K^{+}$ surprise measured against the future. The two axes used
throughout are
\[
  \underbrace{A \;=\; \tfrac{1}{2}\left(K^{-} + K^{+}\right)}_{\text{atypicality}},
  \qquad
  \underbrace{L \;=\; K^{-} - K^{+}}_{\text{lead}} .
\]
$A$ is how unusual the wording is overall, in both temporal directions at once.
$L$ is signed: positive for a leading filing, negative for a lagging one. The
symbol \dkl{} is used interchangeably with $L$ where it aids comparison with
the source literature.

\subsection{Why these two, and not the raw pair}

$A$ and $L$ are a rotation of $(K^{-}, K^{+})$ through 45 degrees. Plotted with
$K^{-}$ on one axis and $K^{+}$ on the other, $A$ is position \emph{along} the
diagonal and $L$ is signed distance \emph{from} it: magnitude and direction of
the same displacement. Three things follow.

The rotation is the natural coordinate system because the raw pair is almost
collinear. $K^{-}$ and $K^{+}$ correlate at $0.979$ on the complete
software-class corpus, so entering both is close to entering one variable
twice; the informative content is the long axis and the small residual, which
is exactly what $A$ and $L$ isolate.

The rotation also very nearly decorrelates the two regressors, which matters
because the paper enters both and interprets them separately. Against $L$, the
correlation of $A$ is $+0.038$; of $K^{-}$ alone, $+0.140$; of $|L|$, $+0.250$
($n = 1{,}107{,}807$). Using either level alone, or the unsigned magnitude, as
the ``how unusual'' axis leaves the two coefficients trading against each
other. Nothing empirical rides on the choice --- given the $0.979$ correlation,
$A$ and $K^{-}$ are nearly the same variable --- but the interpretation is
cleaner and the coefficients are closer to independent.

Finally, the pair is exhaustive in the sense that matters: $(A, L)$ recovers
$(K^{-}, K^{+})$ exactly, so nothing is discarded by working in the rotated
frame.

\begin{table}[h]\centering\small
\caption{The quantities, their sources, and the terms used here. \emph{Lagging}
and \emph{leading} describe the two signs of $L$; \emph{atypicality} is
unsigned and says nothing about direction.}
\label{tab:constructs}
\begin{tabular}{llll}
\toprule
Quantity & Definition & \citeauthor{Murdock2017} & Term used here \\
\midrule
Surprise against the past   & $K^{-}$                        & novelty    & atypicality at filing \\
Surprise against the future & $K^{+}$                        & transience & atypicality in hindsight \\
Their average               & $A = \tfrac{1}{2}(K^{-}+K^{+})$ & ---        & \textbf{atypicality} \\
Their signed difference     & $L = K^{-}-K^{+}$              & resonance  & \textbf{lead} (leading / lagging) \\
Its magnitude               & $|L|$                          & ---        & lead magnitude \\
\bottomrule
\end{tabular}
\end{table}

\subsection{A note on the words}

\emph{Atypicality} is used in a purely textual sense and is not trademark law's
distinctiveness doctrine, which governs registrability along the
generic-to-fanciful spectrum; the two attach to different objects, one to the
mark and one to the description of the goods it covers.
Section~\ref{sec:limits} separates the construct from its other near
neighbours --- novelty, transience, resonance, innovation and radicalness ---
and states in each case what adopting that term would claim.

\emph{Leading} and \emph{lagging} were chosen over the alternatives for the
same reason the causal reading was declined above. To say a filing led its
industry's vocabulary is to place it relative to a trajectory without saying it
produced one, in the way that being ahead of a curve implies the curve without
implying you drew it.
 from
% ssrn_diffusion_paper.tex so it can be revised independently of the rest.
%
% Depends on the host preamble for: \dkl{} (defined as $\Delta$KL), natbib,
% booktabs, amsmath. Citation keys used here all exist in the host
% bibliography: Murdock2017, Barron2018, vonGraevenitz2022.
%
% Naming decided 2026-08-06 after consultation. The signed axis is LEAD, and
% filings are LEADING or LAGGING; "currency" was rejected because readers
% import monetary logic into what is a claim about lexical surprise, and
% "resonance" because it names a mechanism this design cannot observe.
% ---------------------------------------------------------------------------

\section{The construct: how surprising a description is, and when}\label{sec:construct}

\subsection{Two questions asked of every filing}

The measure asks one thing twice. Take a filing's goods/services description,
and group every filing by its Nice class --- one of the 45 international goods
and services categories --- and by its year. Then ask:

\begin{enumerate}[noitemsep]
  \item How surprising would this wording have been to a reader who knew only
        what that industry had filed over the previous five years?
  \item How surprising does the same wording look to a reader who also knows
        what the industry filed over the following five years?
\end{enumerate}

\noindent The two answers are rarely the same, and the gap between them is
informative. Wording that was surprising when filed and unremarkable five
years later was surprising because it came first: the industry moved toward
it. Wording that read as ordinary when filed and unusual afterwards was behind
its industry rather than ahead of it, and was left behind as the field moved
on. Filings of the first kind are \emph{leading}; of the second,
\emph{lagging}.

Both quantities are computable for every filing in the corpus, including the
millions from firms that never patent, and neither requires knowing anything
about the firm.

\subsection{Antecedents, and what does not carry over}

The apparatus is that of \citet{Murdock2017}, developed on Darwin's reading
notebooks, and \citet{Barron2018}, who extended it to French Revolution
parliamentary debates. They score each document against the distribution of
documents before it and again against the distribution after it, and call the
two quantities \emph{novelty} and \emph{transience} and their signed difference
\emph{resonance}.

The arithmetic transfers exactly. The interpretation does not, in two ways that
shape everything downstream.

First, there is no single authoring agent. In a reading notebook the author of
the document is the person whose novelty is being measured, and the reference
stream is that author's own intellectual environment; a high value therefore
says something about a mind. Here each filing is scored against the stream of
\emph{other} firms' filings, and the author is typically counsel drafting for
scope of protection. The same number can only say that a description sits far
from what its industry was writing.

Second, \emph{resonance} names a mechanism this design cannot observe. A filing
can sit where its class was heading because it was imitated, because it and the
class responded to the same external development, or by coincidence, and
nothing here separates those. The vocabulary used below is therefore positional
rather than causal: a leading filing was ahead of its industry's language,
which implies the existence of a trajectory without asserting that the filing
caused it.

The nearest measurement precedent on this corpus is \citet{vonGraevenitz2022},
who flag goods/services tokens that are both new and fast-spreading and use
their geography to study diffusion. Their unit is the token and their novelty
flag is binary, assigned once a token's path is known; the measure here is
continuous and attaches to the filing, so every application carries a position
and that position can be crossed with the mark's own prosecution outcomes.

\subsection{Definitions}

Let $P_t$ be the topic distribution of a filing made in year $t$, and let
$Q_{\mathrm{past}}$ and $Q_{\mathrm{future}}$ aggregate same-class filings over
the $W = 5$ years before and after it. Write
\[
  K^{-} \;=\; \mathrm{KL}\!\left(P_t \,\|\, Q_{\mathrm{past}}\right),
  \qquad
  K^{+} \;=\; \mathrm{KL}\!\left(P_t \,\|\, Q_{\mathrm{future}}\right).
\]
Kullback--Leibler divergence is large when the filing puts weight where its
industry's filings put little, so $K^{-}$ is surprise measured against the
past and $K^{+}$ surprise measured against the future. The two axes used
throughout are
\[
  \underbrace{A \;=\; \tfrac{1}{2}\left(K^{-} + K^{+}\right)}_{\text{atypicality}},
  \qquad
  \underbrace{L \;=\; K^{-} - K^{+}}_{\text{lead}} .
\]
$A$ is how unusual the wording is overall, in both temporal directions at once.
$L$ is signed: positive for a leading filing, negative for a lagging one. The
symbol \dkl{} is used interchangeably with $L$ where it aids comparison with
the source literature.

\subsection{Why these two, and not the raw pair}

$A$ and $L$ are a rotation of $(K^{-}, K^{+})$ through 45 degrees. Plotted with
$K^{-}$ on one axis and $K^{+}$ on the other, $A$ is position \emph{along} the
diagonal and $L$ is signed distance \emph{from} it: magnitude and direction of
the same displacement. Three things follow.

The rotation is the natural coordinate system because the raw pair is almost
collinear. $K^{-}$ and $K^{+}$ correlate at $0.979$ on the complete
software-class corpus, so entering both is close to entering one variable
twice; the informative content is the long axis and the small residual, which
is exactly what $A$ and $L$ isolate.

The rotation also very nearly decorrelates the two regressors, which matters
because the paper enters both and interprets them separately. Against $L$, the
correlation of $A$ is $+0.038$; of $K^{-}$ alone, $+0.140$; of $|L|$, $+0.250$
($n = 1{,}107{,}807$). Using either level alone, or the unsigned magnitude, as
the ``how unusual'' axis leaves the two coefficients trading against each
other. Nothing empirical rides on the choice --- given the $0.979$ correlation,
$A$ and $K^{-}$ are nearly the same variable --- but the interpretation is
cleaner and the coefficients are closer to independent.

Finally, the pair is exhaustive in the sense that matters: $(A, L)$ recovers
$(K^{-}, K^{+})$ exactly, so nothing is discarded by working in the rotated
frame.

\begin{table}[h]\centering\small
\caption{The quantities, their sources, and the terms used here. \emph{Lagging}
and \emph{leading} describe the two signs of $L$; \emph{atypicality} is
unsigned and says nothing about direction.}
\label{tab:constructs}
\begin{tabular}{llll}
\toprule
Quantity & Definition & \citeauthor{Murdock2017} & Term used here \\
\midrule
Surprise against the past   & $K^{-}$                        & novelty    & atypicality at filing \\
Surprise against the future & $K^{+}$                        & transience & atypicality in hindsight \\
Their average               & $A = \tfrac{1}{2}(K^{-}+K^{+})$ & ---        & \textbf{atypicality} \\
Their signed difference     & $L = K^{-}-K^{+}$              & resonance  & \textbf{lead} (leading / lagging) \\
Its magnitude               & $|L|$                          & ---        & lead magnitude \\
\bottomrule
\end{tabular}
\end{table}

\subsection{A note on the words}

\emph{Atypicality} is used in a purely textual sense and is not trademark law's
distinctiveness doctrine, which governs registrability along the
generic-to-fanciful spectrum; the two attach to different objects, one to the
mark and one to the description of the goods it covers.
Section~\ref{sec:limits} separates the construct from its other near
neighbours --- novelty, transience, resonance, innovation and radicalness ---
and states in each case what adopting that term would claim.

\emph{Leading} and \emph{lagging} were chosen over the alternatives for the
same reason the causal reading was declined above. To say a filing led its
industry's vocabulary is to place it relative to a trajectory without saying it
produced one, in the way that being ahead of a curve implies the curve without
implying you drew it.
 from
% ssrn_diffusion_paper.tex so it can be revised independently of the rest.
%
% Depends on the host preamble for: \dkl{} (defined as $\Delta$KL), natbib,
% booktabs, amsmath. Citation keys used here all exist in the host
% bibliography: Murdock2017, Barron2018, vonGraevenitz2022.
%
% Naming decided 2026-08-06 after consultation. The signed axis is LEAD, and
% filings are LEADING or LAGGING; "currency" was rejected because readers
% import monetary logic into what is a claim about lexical surprise, and
% "resonance" because it names a mechanism this design cannot observe.
% ---------------------------------------------------------------------------

\section{The construct: how surprising a description is, and when}\label{sec:construct}

\subsection{Two questions asked of every filing}

The measure asks one thing twice. Take a filing's goods/services description,
and group every filing by its Nice class --- one of the 45 international goods
and services categories --- and by its year. Then ask:

\begin{enumerate}[noitemsep]
  \item How surprising would this wording have been to a reader who knew only
        what that industry had filed over the previous five years?
  \item How surprising does the same wording look to a reader who also knows
        what the industry filed over the following five years?
\end{enumerate}

\noindent The two answers are rarely the same, and the gap between them is
informative. Wording that was surprising when filed and unremarkable five
years later was surprising because it came first: the industry moved toward
it. Wording that read as ordinary when filed and unusual afterwards was behind
its industry rather than ahead of it, and was left behind as the field moved
on. Filings of the first kind are \emph{leading}; of the second,
\emph{lagging}.

Both quantities are computable for every filing in the corpus, including the
millions from firms that never patent, and neither requires knowing anything
about the firm.

\subsection{Antecedents, and what does not carry over}

The apparatus is that of \citet{Murdock2017}, developed on Darwin's reading
notebooks, and \citet{Barron2018}, who extended it to French Revolution
parliamentary debates. They score each document against the distribution of
documents before it and again against the distribution after it, and call the
two quantities \emph{novelty} and \emph{transience} and their signed difference
\emph{resonance}.

The arithmetic transfers exactly. The interpretation does not, in two ways that
shape everything downstream.

First, there is no single authoring agent. In a reading notebook the author of
the document is the person whose novelty is being measured, and the reference
stream is that author's own intellectual environment; a high value therefore
says something about a mind. Here each filing is scored against the stream of
\emph{other} firms' filings, and the author is typically counsel drafting for
scope of protection. The same number can only say that a description sits far
from what its industry was writing.

Second, \emph{resonance} names a mechanism this design cannot observe. A filing
can sit where its class was heading because it was imitated, because it and the
class responded to the same external development, or by coincidence, and
nothing here separates those. The vocabulary used below is therefore positional
rather than causal: a leading filing was ahead of its industry's language,
which implies the existence of a trajectory without asserting that the filing
caused it.

The nearest measurement precedent on this corpus is \citet{vonGraevenitz2022},
who flag goods/services tokens that are both new and fast-spreading and use
their geography to study diffusion. Their unit is the token and their novelty
flag is binary, assigned once a token's path is known; the measure here is
continuous and attaches to the filing, so every application carries a position
and that position can be crossed with the mark's own prosecution outcomes.

\subsection{Definitions}

Let $P_t$ be the topic distribution of a filing made in year $t$, and let
$Q_{\mathrm{past}}$ and $Q_{\mathrm{future}}$ aggregate same-class filings over
the $W = 5$ years before and after it. Write
\[
  K^{-} \;=\; \mathrm{KL}\!\left(P_t \,\|\, Q_{\mathrm{past}}\right),
  \qquad
  K^{+} \;=\; \mathrm{KL}\!\left(P_t \,\|\, Q_{\mathrm{future}}\right).
\]
Kullback--Leibler divergence is large when the filing puts weight where its
industry's filings put little, so $K^{-}$ is surprise measured against the
past and $K^{+}$ surprise measured against the future. The two axes used
throughout are
\[
  \underbrace{A \;=\; \tfrac{1}{2}\left(K^{-} + K^{+}\right)}_{\text{atypicality}},
  \qquad
  \underbrace{L \;=\; K^{-} - K^{+}}_{\text{lead}} .
\]
$A$ is how unusual the wording is overall, in both temporal directions at once.
$L$ is signed: positive for a leading filing, negative for a lagging one. The
symbol \dkl{} is used interchangeably with $L$ where it aids comparison with
the source literature.

\subsection{Why these two, and not the raw pair}

$A$ and $L$ are a rotation of $(K^{-}, K^{+})$ through 45 degrees. Plotted with
$K^{-}$ on one axis and $K^{+}$ on the other, $A$ is position \emph{along} the
diagonal and $L$ is signed distance \emph{from} it: magnitude and direction of
the same displacement. Three things follow.

The rotation is the natural coordinate system because the raw pair is almost
collinear. $K^{-}$ and $K^{+}$ correlate at $0.979$ on the complete
software-class corpus, so entering both is close to entering one variable
twice; the informative content is the long axis and the small residual, which
is exactly what $A$ and $L$ isolate.

The rotation also very nearly decorrelates the two regressors, which matters
because the paper enters both and interprets them separately. Against $L$, the
correlation of $A$ is $+0.038$; of $K^{-}$ alone, $+0.140$; of $|L|$, $+0.250$
($n = 1{,}107{,}807$). Using either level alone, or the unsigned magnitude, as
the ``how unusual'' axis leaves the two coefficients trading against each
other. Nothing empirical rides on the choice --- given the $0.979$ correlation,
$A$ and $K^{-}$ are nearly the same variable --- but the interpretation is
cleaner and the coefficients are closer to independent.

Finally, the pair is exhaustive in the sense that matters: $(A, L)$ recovers
$(K^{-}, K^{+})$ exactly, so nothing is discarded by working in the rotated
frame.

\begin{table}[h]\centering\small
\caption{The quantities, their sources, and the terms used here. \emph{Lagging}
and \emph{leading} describe the two signs of $L$; \emph{atypicality} is
unsigned and says nothing about direction.}
\label{tab:constructs}
\begin{tabular}{llll}
\toprule
Quantity & Definition & \citeauthor{Murdock2017} & Term used here \\
\midrule
Surprise against the past   & $K^{-}$                        & novelty    & atypicality at filing \\
Surprise against the future & $K^{+}$                        & transience & atypicality in hindsight \\
Their average               & $A = \tfrac{1}{2}(K^{-}+K^{+})$ & ---        & \textbf{atypicality} \\
Their signed difference     & $L = K^{-}-K^{+}$              & resonance  & \textbf{lead} (leading / lagging) \\
Its magnitude               & $|L|$                          & ---        & lead magnitude \\
\bottomrule
\end{tabular}
\end{table}

\subsection{A note on the words}

\emph{Atypicality} is used in a purely textual sense and is not trademark law's
distinctiveness doctrine, which governs registrability along the
generic-to-fanciful spectrum; the two attach to different objects, one to the
mark and one to the description of the goods it covers.
Section~\ref{sec:limits} separates the construct from its other near
neighbours --- novelty, transience, resonance, innovation and radicalness ---
and states in each case what adopting that term would claim.

\emph{Leading} and \emph{lagging} were chosen over the alternatives for the
same reason the causal reading was declined above. To say a filing led its
industry's vocabulary is to place it relative to a trajectory without saying it
produced one, in the way that being ahead of a curve implies the curve without
implying you drew it.
 from
% ssrn_diffusion_paper.tex so it can be revised independently of the rest.
%
% Depends on the host preamble for: \dkl{} (defined as $\Delta$KL), natbib,
% booktabs, amsmath. Citation keys used here all exist in the host
% bibliography: Murdock2017, Barron2018, vonGraevenitz2022.
%
% Naming decided 2026-08-06 after consultation. The signed axis is LEAD, and
% filings are LEADING or LAGGING; "currency" was rejected because readers
% import monetary logic into what is a claim about lexical surprise, and
% "resonance" because it names a mechanism this design cannot observe.
% ---------------------------------------------------------------------------

\section{The construct: how surprising a description is, and when}\label{sec:construct}

\subsection{Two questions asked of every filing}

The measure asks one thing twice. Take a filing's goods/services description,
and group every filing by its Nice class --- one of the 45 international goods
and services categories --- and by its year. Then ask:

\begin{enumerate}[noitemsep]
  \item How surprising would this wording have been to a reader who knew only
        what that industry had filed over the previous five years?
  \item How surprising does the same wording look to a reader who also knows
        what the industry filed over the following five years?
\end{enumerate}

\noindent The two answers are rarely the same, and the gap between them is
informative. Wording that was surprising when filed and unremarkable five
years later was surprising because it came first: the industry moved toward
it. Wording that read as ordinary when filed and unusual afterwards was behind
its industry rather than ahead of it, and was left behind as the field moved
on. Filings of the first kind are \emph{leading}; of the second,
\emph{lagging}.

Both quantities are computable for every filing in the corpus, including the
millions from firms that never patent, and neither requires knowing anything
about the firm.

\subsection{Antecedents, and what does not carry over}

The apparatus is that of \citet{Murdock2017}, developed on Darwin's reading
notebooks, and \citet{Barron2018}, who extended it to French Revolution
parliamentary debates. They score each document against the distribution of
documents before it and again against the distribution after it, and call the
two quantities \emph{novelty} and \emph{transience} and their signed difference
\emph{resonance}.

The arithmetic transfers exactly. The interpretation does not, in two ways that
shape everything downstream.

First, there is no single authoring agent. In a reading notebook the author of
the document is the person whose novelty is being measured, and the reference
stream is that author's own intellectual environment; a high value therefore
says something about a mind. Here each filing is scored against the stream of
\emph{other} firms' filings, and the author is typically counsel drafting for
scope of protection. The same number can only say that a description sits far
from what its industry was writing.

Second, \emph{resonance} names a mechanism this design cannot observe. A filing
can sit where its class was heading because it was imitated, because it and the
class responded to the same external development, or by coincidence, and
nothing here separates those. The vocabulary used below is therefore positional
rather than causal: a leading filing was ahead of its industry's language,
which implies the existence of a trajectory without asserting that the filing
caused it.

The nearest measurement precedent on this corpus is \citet{vonGraevenitz2022},
who flag goods/services tokens that are both new and fast-spreading and use
their geography to study diffusion. Their unit is the token and their novelty
flag is binary, assigned once a token's path is known; the measure here is
continuous and attaches to the filing, so every application carries a position
and that position can be crossed with the mark's own prosecution outcomes.

\subsection{Definitions}

Let $P_t$ be the topic distribution of a filing made in year $t$, and let
$Q_{\mathrm{past}}$ and $Q_{\mathrm{future}}$ aggregate same-class filings over
the $W = 5$ years before and after it. Write
\[
  K^{-} \;=\; \mathrm{KL}\!\left(P_t \,\|\, Q_{\mathrm{past}}\right),
  \qquad
  K^{+} \;=\; \mathrm{KL}\!\left(P_t \,\|\, Q_{\mathrm{future}}\right).
\]
Kullback--Leibler divergence is large when the filing puts weight where its
industry's filings put little, so $K^{-}$ is surprise measured against the
past and $K^{+}$ surprise measured against the future. The two axes used
throughout are
\[
  \underbrace{A \;=\; \tfrac{1}{2}\left(K^{-} + K^{+}\right)}_{\text{atypicality}},
  \qquad
  \underbrace{L \;=\; K^{-} - K^{+}}_{\text{lead}} .
\]
$A$ is how unusual the wording is overall, in both temporal directions at once.
$L$ is signed: positive for a leading filing, negative for a lagging one. The
symbol \dkl{} is used interchangeably with $L$ where it aids comparison with
the source literature.

\subsection{Why these two, and not the raw pair}

$A$ and $L$ are a rotation of $(K^{-}, K^{+})$ through 45 degrees. Plotted with
$K^{-}$ on one axis and $K^{+}$ on the other, $A$ is position \emph{along} the
diagonal and $L$ is signed distance \emph{from} it: magnitude and direction of
the same displacement. Three things follow.

The rotation is the natural coordinate system because the raw pair is almost
collinear. $K^{-}$ and $K^{+}$ correlate at $0.979$ on the complete
software-class corpus, so entering both is close to entering one variable
twice; the informative content is the long axis and the small residual, which
is exactly what $A$ and $L$ isolate.

The rotation also very nearly decorrelates the two regressors, which matters
because the paper enters both and interprets them separately. Against $L$, the
correlation of $A$ is $+0.038$; of $K^{-}$ alone, $+0.140$; of $|L|$, $+0.250$
($n = 1{,}107{,}807$). Using either level alone, or the unsigned magnitude, as
the ``how unusual'' axis leaves the two coefficients trading against each
other. Nothing empirical rides on the choice --- given the $0.979$ correlation,
$A$ and $K^{-}$ are nearly the same variable --- but the interpretation is
cleaner and the coefficients are closer to independent.

Finally, the pair is exhaustive in the sense that matters: $(A, L)$ recovers
$(K^{-}, K^{+})$ exactly, so nothing is discarded by working in the rotated
frame.

\begin{table}[h]\centering\small
\caption{The quantities, their sources, and the terms used here. \emph{Lagging}
and \emph{leading} describe the two signs of $L$; \emph{atypicality} is
unsigned and says nothing about direction.}
\label{tab:constructs}
\begin{tabular}{llll}
\toprule
Quantity & Definition & \citeauthor{Murdock2017} & Term used here \\
\midrule
Surprise against the past   & $K^{-}$                        & novelty    & atypicality at filing \\
Surprise against the future & $K^{+}$                        & transience & atypicality in hindsight \\
Their average               & $A = \tfrac{1}{2}(K^{-}+K^{+})$ & ---        & \textbf{atypicality} \\
Their signed difference     & $L = K^{-}-K^{+}$              & resonance  & \textbf{lead} (leading / lagging) \\
Its magnitude               & $|L|$                          & ---        & lead magnitude \\
\bottomrule
\end{tabular}
\end{table}

\subsection{A note on the words}

\emph{Atypicality} is used in a purely textual sense and is not trademark law's
distinctiveness doctrine, which governs registrability along the
generic-to-fanciful spectrum; the two attach to different objects, one to the
mark and one to the description of the goods it covers.
Section~\ref{sec:limits} separates the construct from its other near
neighbours --- novelty, transience, resonance, innovation and radicalness ---
and states in each case what adopting that term would claim.

\emph{Leading} and \emph{lagging} were chosen over the alternatives for the
same reason the causal reading was declined above. To say a filing led its
industry's vocabulary is to place it relative to a trajectory without saying it
produced one, in the way that being ahead of a curve implies the curve without
implying you drew it.
